% Why the release schedule alone cannot matter --- plan \S3.3, "\S2".
% Sources: 03_fixed_removal.tex:5, :7, :12-15, :16-17; figure from
% 01_opening.tex:84-120.
%
% PRESERVED NEAR-VERBATIM, do not restyle (plan \S3.6, report item 2):
%   * the negative result, 03:5-9;
%   * the Brownian-drift and macrophage-motility caveats, 03:16-17, which
%     \ChDist{} depends on (ch5_handoffs.md entry 3).

\section{Why the release schedule alone cannot matter}
\label{path:sec:schedule}

In the pursuit of modelling multi-stage intracellular infection, an obvious
first move is an iterated birth--death--catastrophe process with the
assumption of complete transfer --- chained ruptures. In trying to determine
whether there is a fundamental reason for bursting dynamics in pathogenesis,
this model is of limited efficacy. With agents replicating identically
whenever they are within a host cell, it matters not how they are moved,
just that they are replicating. With chained ruptures, the burst event is
merely a recorded instance without tangible implication to agent spread. And
if instead the bacteria modelled were to make their way one by one to a new
host, or even to multiple hosts, no change would be made to the expected
growth rate of the overall population. Even if some delay were introduced
for unit transfer between hosts, this would merely have the effect of adding
intervals of inactivity to the time axis; broad propagation dynamics would
remain identical.

It seems, then, that if we wish to understand why pathogens follow a
bursting regime of infection --- or whether there is a reason at all --- we
must consider additional factors. Namely, we must take into account that the
death and replication rates of individual agents are not entirely
independent of the surrounding cohort. Various of the body's immune factors
are depletable. Leucocytes including macrophages and neutrophils contain on
the order of ten so-called granules, small packets containing oxidising
compounds with a toxic effect on invading pathogens. Given that these can be
used up, the effective within-host death rate of infectious agent will
reduce as more and more of the incident cohort are killed. There are likely
extracellular immune factors which exhibit a similar pattern of depletion.

\subsection{Successive birth--death--catastrophe}
\label{path:sec:chainedruptures}

The concept of chained ruptures is simple. A series of
birth--death--catastrophe processes follow on from one another, the result
of the last going on to found an initial condition for the next
\cite{brockwell1986extinction,stirzaker2006processes,di2008note}. Starting
with a single replicative unit, the first process undergoes the standard
dynamics of birth and death. Then, at the time of the catastrophe event, the
quantity of units present immediately transfers into another target cell,
founding the second process with its initial condition as given. This is the
main assumption: that released replicative units somehow find their way, all
at once, to the same host cell. Of course this is quite reduced from the
biological context. \ChDistcap{} solves the chain in closed form at $\mu=0$ for
every chain position, giving the rupture-size law $\pr{r_k=n} =
\binom{n+k-2}{k-1}s^k\rho^{\,n-1}$ of \cref{dist:eq:rklaw} together with its
moments and rupture times, and records the case $\mu>0$ as open
(\cref{dist:sec:chained-scope}).

A less restricted view would be that following rupture, all of the $r_1$
released bacteria undergo Brownian motion and drift gradually from the burst
location. The cohort, as it spreads out, will encounter multiple new target
cells and get taken up one by one. Albeit less so, this is also greatly
reduced from biological reality, which in its full complexity will be beyond
the scope of concise mathematical description. For example, the target cells
we have focused on mostly are macrophages. Macrophages do not merely float
around in Brownian drift; they are motile, and essentially walk around their
environment seeking out debris and alien material to engulf and make away
with. The agents themselves are little more straightforward. \yp{} can
``decide'' whether it will be absorbed or not by such phagocytes
\cite{MonackEtAl1997,PujolBliska2005}. In addition to this, the bacteria can
inject immune cells with a host of special proteins with various effects
including neutralisation \cite{KeEtAl2013}.

Generally, the playing field of these actors is just that, an environment of
play, in which the responsive agents interact with one another; causal
relationships ensue which go far beyond mere drift and chance.

\begin{figure}[htbp]
    \centering
    \setlength{\tabcolsep}{1.5pt}
    \begin{tabular}{@{}cccc@{}}
      \multicolumn{4}{c}{\footnotesize time $\longrightarrow$} \\[1pt]
      \includegraphics[height=0.23\textwidth]{basic1.png}
      & \includegraphics[height=0.23\textwidth]{basic2.png}
      & \includegraphics[height=0.23\textwidth]{basic3.png}
      & \includegraphics[height=0.23\textwidth]{basic4.png} \\
      {\footnotesize (a)} & {\footnotesize (b)}
      & {\footnotesize (c)} & {\footnotesize (d)} \\
    \end{tabular}
    \caption{Four frames of a particle simulation, showing what the chained
    model leaves out. (a) Intracellular replication before bursting: the
    load accumulates out of reach. (b) The cell bursts and the contents
    span the extracellular space before finding new hosts --- the interval
    the chained model collapses to a point. (c) Macrophages are motile;
    they wander the extracellular matrix phagocytosing debris and
    extracellular pathogens alike, so which cell a released agent reaches
    is not settled by diffusion alone. (d) A second round of bursting, here
    with a larger load. Under identical intracellular rates the growth rate
    is the same however the transfers between (a) and (d) are arranged,
    which is why the argument of this section forces a nonlinearity into
    the model.}
    \label{path:fig:basic}
\end{figure}

\Cref{path:fig:basic} shows the four stages of that picture as a simulation
run. What matters for the argument is the contrast between panels (a) and
(d): both show a burst, and under identical intracellular rates the two
carry the same expected yield whatever route the released agents take
between them.

The conclusion is negative, and it sets the chapter's agenda. A release
schedule can only matter if something in the environment responds to how
many agents arrive at once --- if the per-agent hazard depends on the size
of the cohort it arrives in. Two such mechanisms are treated in what
follows, one inside the host cell and one outside it:
\cref{path:sec:intracellular} makes the intracellular killing capacity
finite and depletable, and \cref{path:sec:removal} does the same for the
extracellular one. Both are nonlinear, and both are therefore outside the
scope of the branching comparison in \ChPop{}, which compares offspring
distributions at fixed mean and says so explicitly
(\cref{p:rem:flood-scope}).

\FloatBarrier
